% ID: SEL-D0018-BB-E04
% ESTADO: cotejada. Segunda pasada visual completa y borrador compilado revisados.
% Fuente: Selectividad/matematicas_aplicadas/matematicas_aplicadas_P_F.pdf
% SHA-256: fe6fd5b48407569d41169c80d71e6997050d6d06d1afd41cd639a553f9b712b5
% Página: 1; bloque B; ejercicio original 4.

Se considera la función
\[
f(x)=\begin{cases}
\dfrac{2}{x+1} & \text{si }x<-2\\[6pt]
x^2+a & \text{si }x\geq -2
\end{cases}
\]
\begin{enumerate}[label=\alph*)]
\item \textbf{(1 punto)} Calcule el valor de $a$ para que $f$ sea continua en todo su dominio. Para ese valor de $a$, \textquestiondown{}es derivable la función $f$?
\item \textbf{(0.5 puntos)} Para $a=-6$, halle la ecuación de la recta tangente a la gráfica de $f$ en el punto de abscisa $x=3$.
\item \textbf{(1 punto)} Para $a=-6$, esboce la gráfica de $f$ y calcule el área de la región limitada por la gráfica de la función $f$, el eje de abscisas y las rectas $x=3$ y $x=5$.
\end{enumerate}
